% Fragmento público generado desde propietarios íntegros.
% Residencia: 52.2.2.
% Fuente: ../incoming/radion_monografia_20260827/manuscrito/sections/03_cierre_exacto.tex:51-103

\paragraph{Oriented density \(20/81\) of the active sector \(90/120\)}

The triadic phase separator selects
\[
H_+=\{1,4,7\},\qquad H_-=\{2,5,8\},\qquad H_0=\{3,6,9\}.
\]
There are six active phases and fifteen unordered pairs:
\[
H_{\mathrm{act}}=H_+\cup H_-,\qquad
\left|\binom{H_{\mathrm{act}}}2\right|=\binom62=15.
\]
The active mode and its extension to eight positions are
\begin{equation}
N_6=6\binom62=90,\qquad
N_8=8\binom62=120.
\label{eq:90-120-incidencia}
\end{equation}
Two conjugate sheets of the active sector, normalized by the ambient fiber,
give
\begin{equation}
\rho_{\mathrm{tw}}
=\frac{2N_6}{|\F_3^6|}
=\frac{2\cdot90}{729}
=\frac{180}{729}
=\boxed{\frac{20}{81}}.
\label{eq:rho-twoway}
\end{equation}

\begin{proposition}[Relative uniqueness of the density]
With the active sector \(N_6=90\), the two conjugate sheets, and the
normalization by \(|\F_3^6|=729\) fixed, the oriented density is necessarily
\(20/81\).
\end{proposition}
\begin{proof}
The transported cardinality is \(2N_6=180\). Normalization by the ambient space
gives \(180/729\); dividing the numerator and denominator by \(9\) yields \(20/81\).
No comparison with \(180/\pi\) or with \(\epsR\) is involved.
\end{proof}

The historical reading as a projection onto a \(1/81\) lattice is a useful rigidity control, but it is no longer the foundation: the coefficient is obtained before closure. Ablating one sheet produces \(10/81\), and the output changes at order \(10^{-5}\); hence, the factor of two is not decorative.

The torsional section is
\begin{equation}
\Phi_T=\rho_{\mathrm{tw}}\Delta_4,
\qquad
S_T=1+\Phi_T.
\label{eq:seccion-torsional}
\end{equation}

